\documentclass[12pt]{article}    %-*-latex-*-
%
\font\fff=eufm10 scaled \magstep1
%\font\STM=stmary10
\font\STM=cmss10
\parindent=0pt
%
%\def\g{\hbox{\fff g}}
%
\begingroup 
\newtheorem{bigthm}{Theorem}             % Numbered separately, as A, B, etc.
\newtheorem{thm}{Theorem}[section]       % Numbered within each section
\newtheorem{lemma}[thm]{Lemma}
\newtheorem{proposition}[thm]{Proposition}
\newtheorem{theorem}[thm]{Theorem}
\newtheorem{corollary}[thm]{Corollary}
\newtheorem{definition}[thm]{Definition}
\newtheorem{notation}{Notation}
\newtheorem{terminology}{Terminology}
\newtheorem{remark}[thm]{Remark}          % Numbered along with thm
\newtheorem{example}[thm]{Example}        % Numbered along with thm
\newtheorem{observation}[thm]{Observation}        % Numbered along with thm
\endgroup
%
%dirac:
\def\dir#1.{\;|\, #1 \,\rangle}             % dirac bracket
\def\dirs#1.{\;|\, #1 \,\rangle_{\rm Sym}}  % dirac symmetric bracket
%\def\state#1.{|#1\rangle}
%\def\.#1.{\langle #1 \rangle_{wv}}
%
\def\qed{\qquad\framebox[7pt]\medskip\noindent}
%\def\pqed{\rightline{\qed}\medskip\noindent}
%
\def\FFF#1#2{\Phi^{(#1)}_{#2}}
\def\tr{\mathrm{tr}\,}
%\def\.#1.#2.{{#1}^{#2}}
%\def\tfo #1.#2.#3.#4.{{}_2 F_1\left({#1,#2\atop #3}\biggm| #4 \right)}
%\def\kck#1.#2.{{#1\choose#2}}
%\def\hft #1.#2.#3.#4.#5.#6.{{}_3 F_2\left({#1,\,#2,\,#3\atop#4,\,#5}\biggm| #6 \right)}
%\def\ov#1.#2.{\frac{#1}{#2}}
%
%
\newenvironment{proof}{{\sl Proof:}\quad}{\hfill{\qed}\\ \noindent}
%
\renewcommand{\thebigthm}{\Alph{bigthm}}  % Number as "Theorem A."
\renewcommand{\thenotation}{}  
\renewcommand{\theterminology}{}  
%
\renewenvironment{abstract}{\begin{quote}{\bf Abstract.\
}\small}{\end{quote}\bigskip}
%
%
%
\title{Krawtchouk matrices from Hadamard matrices}
\author{Philip Feinsilver and Jerzy Kocik
\thanks{\leftskip -1cm Department of Mathematics
Southern Illinois University Carbondale, IL. 62901, U.S.A.} 
} 

\date{}

\thispagestyle{empty}

\begin{document} 

\maketitle

\begin{abstract}
We show the connection between the Krawtchouk matrices (and hence,
polynomials) and the Sylvester-Hadamard matrices. Then we derive the
Krawtchouk matrices from the classical symmetric Bernoulli random
walk. This approach is continued to the quantum probability context, 
leading to the construction of Krawtchouk matrices via tensor powers 
of the elementary Krawtchouk/Hadamard matrix. 

{\bf Keywords:} Krawtchouk polynomials, Hadamard matrices, 
symmetric tensors, random walk, quantum computing, quantum probability.

{\bf AMS classification:} 
05B20, %hadamard
60G50, %random walk
47A80, %tensor, operators 
81P99, %qm computing
46L53. %noncom probab


\end{abstract}


%\vfill
%\pagebreak
%\tableofcontents
%\pagebreak

%===================  1   ===========================================
\section{Introduction}
\label{sec:intro}

Krawtchouk polynomials were introduced in \cite{Kra} and see also \cite{Kra2} 
B: digital filtering
Del,Del2,Del3:association schemes
Dun,DR: harmonic analysis/special functions
Eag,V-J: statistics
FF,FS: symmetric K-matrices, classical probabilistic description
FH: tensor spaces
Lev: coding theory
MS: coding theory
Spr: symm. trace th.
Sze:orthogonal polynomials
Y-H: Hadamard matrices/transform
(see also 
\cite{B,Dun, DR, Eag, Kra2, V-J} and especially \cite{Sze}). They 
play important role in many areas.  They appear in code theory 
(used explicitly first in \cite{Del, Del2, Del3}) where 
MacWilliams transform (see \cite{MS} page 137) relates dual codes. 


%\bibitem{Bose}\bibitem{Del}\bibitem{Del2}\bibitem{Del3}\bibitem{Dun}
%\bibitem{DR}\bibitem{Eag}\bibitem{FF}\bibitem{FS}\bibitem{FH}
%\bibitem{Kra}
%\bibitem{Kra2}
%\bibitem{Lev}\bibitem{MS}\bibitem{Spr}\bibitem{Sze}\bibitem{V-J}
%\bibitem{Wie}\bibitem{YH} 


%%%%%%%%%%%%%%
\vfill       %  To be removed
\newpage     %
%%%%%%%%%%%%%%

%===================  1   ===========================================
\section{What are Krawtchouk  Matrices}

Krawtchouk matrices are generalization of binomial coefficients.
For instance the ``degree-two algebraic rules" 
translate into a ``second-degree Krawtchouk  matrix": 

$$
\matrix{  
       \hfill (a+b)^2   &=& a^2+2ab+b^2         \hfill \cr
       \hfill(a+b)(a-b) &=& a^2-ab+b^2          \hfill \cr
       \hfill (a-b)^2   &=& a^2-3a^2b+3ab^2+b^3 \hfill \cr}
\qquad\Rightarrow\qquad
\Phi^{(2)} = \left[\matrix{  1 &  1 &  1 \cr
                             2 &  0 & -2 \cr
                             1 & -1 &  1 \cr }\right]
$$
(the coefficients are set vertically).
Hence the formal definition:

%----------------------
\begin{definition}\rm
 
The $n^{\mathrm{th}}$-degree Krawtchouk  matrix $\FFF n{}$ is an
$(n+1)\times(n+1)$ matrix,
the entries of which are determined by:
\begin{equation}
\label{eq:genkraw}
    (1+v)^{n-j} \; (1-v)^j = \sum_{i=0}^{n} \ v^i \Phi^{(n)}_{ij}
\end{equation}
where the left-hand-side
$
              G(v)= (1+v)^{n-j}\; (1-v)^j
$
can be viewed as the {\it generating function} for $\Phi$.
Simple exercise gives direct expression:
$$
\Phi^{(n)}_{ij}= \sum_{\alpha} (-)^\alpha {j \choose\alpha}
      {n-j \choose i-\alpha}
$$ 
\end{definition}
%-----------------------------

Krawtchouk  matrices of order zero and one are:
%
\begin{equation}
\label{eq:krav2}
\Phi^{(0)}=\left[\matrix{ 1 }\right] \qquad
\Phi^{(1)}=\left[\matrix{ 1 &  1 \cr
                          1 & -1 \cr}\right]
\end{equation}

The columns of the Krawtchouk matrices contain the generalized 
binomials, while the rows define Krawtchouk polynomials: for the order 
$n$, define the $p$-th Krawtchouk| polynomial as a function that assumes 
the values  in the $i$-th row. 
$$
        P^{(n)}_i(j) = \Phi^{(n)}_{ij}   
$$ 
One can easily show that $P_i(j)$ is indeed a degree $i$ 
polynomial in $j$. 
\\


%-----------------------------------------------------------
The basic properties of Krawtchouk matrices are:
\\

(1)
Square of a Krawtchouk  matrix is proportional to the identity matrix.
$$
\FFF n{} \circ \FFF n {} = 2^n \cdot \hbox{id}
$$
This remarkable property allows one to defines 
a Fourier-like transformation for integer vectors. 
\\

(2)
The leftmost entries are just binomial coefficients,
$\FFF n{i0} = {n\choose i}$.  The rightmost entries are just binomial
coefficients with the alternating sign, $\FFF n{0i} = (-1)^i{n\choose i}$
The top row has always 1's only. The bottom row has $\pm 1$'s with
alternating signs (starting with $+1$).
\\

(3)
There is a four-fold symmetry:
$
\mid \FFF{n}{i\;  j} \mid  =  \mid \FFF{n}{n-i\;j} \mid =
\mid \FFF{n}{i\;n-j} \mid =  \mid \FFF{n}{n-i\;n-j} \mid
$.
\\


If each column of the matrix is multiplied by a corresponding binomial
coefficient, the matrix becomes symmetric.
Let $B^{(n)}$ denote an $(n+1)\times (n+1)$ matrix with binomial
coefficients on its diagonal, 
%
\begin{equation}
\label{eq:diagB}
          B^{(n)}_{ii}=\hbox{\rm diag}{n\choose i}
\end{equation}
and with zeros everywhere else.  
Then, one defines the  {\bf symmetric Krawtchouk  matrices}
as
$$
S^{(n)}=\FFF n{} \circ B^{(n)}
$$
{\bf Example:}  For $n=3$, we have
$$
S^{3)} =
       \left[\matrix{ 1 &  1 &  1  &  1 \cr
                      3 &  1 & -1  & -3 \cr
                      3 & -1 & -1  &  3 \cr
                      1 & -1 &  1  & -1 \cr }\right]
\circ
       \left[\matrix{ 1 &  0 & 0 & 0 \cr
                      0 &  3 & 0 & 0 \cr
                      0 &  0 & 3 & 0 \cr
                      0 &  0 & 0 & 1 \cr }\right]
=
      \left[\matrix{ 1 &  3 &  3  &  1 \cr
                     3 &  3 & -3  & -3 \cr
                     3 & -3 & -3  &  3 \cr
                     1 & -3 &  3  & -1 \cr }\right]
$$ 
A number of symmetric Krawtchouk  matrices are displayed in Table 2.  
\\




%==================  krav matrices ==================================
\newpage
\hrule

\begin{eqnarray*}
%
\Phi^{(0)}&=&\left[\matrix{ 1 }\right]   
\\
\Phi^{(1)}&=&\left[\matrix{ 1 &  1 \cr
                           1 & -1 \cr}\right]    
\\
\Phi^{(2)}&=& \left[\matrix{  1 &  1 &  1 \cr
                             2 &  0 & -2 \cr
                             1 & -1 &  1 \cr }\right]
\\
\Phi^{(3)}&=&
       \left[\matrix{ 1 &  1 &  1  &  1 \cr
                      3 &  1 & -1  & -3 \cr
                      3 & -1 & -1  &  3 \cr
                      1 & -1 &  1  & -1 \cr }\right]
\\
\Phi^{(4)}&=&
       \left[\matrix{ 1 &  1 &  1  &  1  &  1 \cr
                      4 &  2 &  0  & -2  & -4 \cr
                      6 &  0 & -2  &  0  &  6 \cr
                      4 & -2 &  0  &  2  & -4 \cr
                      1 & -1 &  1  & -1  &  1 \cr }\right]
\\
\Phi^{(5)}&=&
       \left[\matrix{ 1  &  1 &  1  &  1  &  1 &   1 \cr
                      5  &  3 &  1  & -1  & -3 &  -5 \cr
                     10  &  2 & -2  & -2  &  2 &  10 \cr
                     10  & -2 & -2  &  2  &  2 & -10 \cr
                      5  & -3 &  1  &  1  & -3 &   5 \cr
                      1  & -1 &  1  & -1  &  1 &  -1 \cr }\right]
\\
\Phi^{(6)}&=&
       \left[\matrix{ 1  &  1 &  1  &  1  &  1 &  1  &   1 \cr
                      6  &  4 &  2  &  0  & -2 & -4  &  -6 \cr
                     15  &  5 & -1  & -3  & -1 &  5  &  15 \cr
                     20  &  0 & -4  &  0  &  4 &  0  & -20 \cr
                     15  & -5 & -1  &  3  & -1 & -5  &  15 \cr
                      6  & -4 &  2  &  0  & -2 &  4  &  -6 \cr
                      1  & -1 &  1  & -1  &  1 & -1  &   1 \cr }\right]
\end{eqnarray*}



\hrule
\bigskip
{\bf Table 1:} Krawtchouk  matrices

%=============   symmetric krav matrices =============================
\newpage
\hrule


\begin{eqnarray*}
%
S^{(0)}&=&\left[\matrix{ 1 }\right]
\\
S^{(1)}&=&\left[\matrix{ 1 &  1 \cr
                        1 & -1 \cr}\right]
\\
S^{(2)}&=& \left[\matrix{  1 &  2 &  1 \cr
                          2 &  0 & -2 \cr
                          1 & -2 &  1 \cr }\right]
\\
S^{(3)}&=&
      \left[\matrix{ 1 &  3 &  3  &  1 \cr
                     3 &  3 & -3  & -3 \cr
                     3 & -3 & -3  &  3 \cr
                     1 & -3 &  3  & -1 \cr }\right]
\\
\\
S^{(4)}&=&
      \left[\matrix{ 1 &  4 &  6  &  4  &  1 \cr
                     4 &  8 &  0  & -8  & -4 \cr
                     6 &  0 &-12  &  0  &  6 \cr
                     4 & -8 &  0  &  8  & -4 \cr
                     1 & -4 &  6  & -4  &  1 \cr }\right]
\\
S^{(5)}&=&
      \left[\matrix{ 1  &  5 & 10 &  10  &   5 &   1 \cr
                     5  & 15 & 10 & -10  & -15 &  -5 \cr
                    10  & 10 &-20 & -20  &  10 &  10 \cr
                    10  &-10 &-20 &  20  &  10 & -10 \cr
                     5  &-15 & 10 &  10  & -15 &   5 \cr
                     1  & -5 & 10 & -10  &   5 &  -1 \cr }\right]
\\
S^{(6)}&=&
       \left[\matrix{ 1  &   6 &  15  &  20 &  15 &   6 &   1 \cr
                      6  &  24 &  30  &   0 & -30 & -24 &  -6 \cr
                     15  &  30 & -15  & -60 & -15 &  30 &  15 \cr
                     20  &   0 & -60  &   0 &  60 &   0 & -20 \cr
                     15  & -30 & -15  &  60 & -15 & -30 &  15 \cr
                      6  & -24 &  30  &   0 & -30 &  24 &  -6 \cr
                      1  &  -6 &  15  & -20 &  15 &  -6 &   1 \cr }\right]
\end{eqnarray*}

\hrule
\bigskip
{\bf Table 2:} Symmetric Krawtchouk  matrices






% ===============================================

\newpage


%===================  2   ===========================================
\section{Krawtchouk tower and identities}
\label{sec:kraw}

Krawtchouk  matrices generalize Pascal's triangle.
One can visualize that in the form of ``Pascal-Krawtchouk pyramid" 
by setting the matrices one upon the other, increasing the degree 
$n$ down-wards
(Pascal triangle forms one side of Krawtchouk tower).
\\

As one could expect, Pascal identity, valid for the leftmost columns, 
generalizes for the other columns Krawtchouk  matrices.
Less obvious is another ---call it dual Pascal--- identity.
\\


\begin{lemma}
{\bf (Cross identities)}
Denote
$$\left\{ \matrix{
  a&=& \FFF n {i-1\;j} \cr
  b&=& \FFF n {i \;j} \cr } \right.
  \qquad\left\{ \matrix{
  A&=& \FFF {n+1} {i\;j} \cr
  B&=& \FFF {n+1} {i\;j+1} \cr}  \right.
$$
Then one has
$$
\matrix{
      (i)&\quad &a+b &=& A &\qquad&(iii) & A+B &=& 2b   \cr
     (ii)&\quad &b-a &=& B &\qquad&(iv)  &  A-B &=& 2a  \cr }
$$
\end{lemma}

\begin{proof}
To prove (i), consider 
$(1+v)^{n-j+1}(1-v)^j=(1+v)\times (1+v)^{n-j}(1-v)^j$.
To prove (ii), consider 
$(1+v)^{n-j}(1-v)^{j+1}= (1-v) \times  (1+v)^{n-j}(1-v)^j$.
Now equations (iii) and (iv) follow from (i) and (ii).
\end{proof}


\begin{corollary}
{\bf (Square identity)}
In a square of any four adjacent entries in a Krawtchouk  matrix, 
the entry in the left-bottom corner is the sum of the other three,
i.e., 
$$
\hbox{for } ~~ 
\Phi =  
\left[\matrix{  &\vdots&\vdots \cr
              \cdots & a & c &\cdots \cr
              \cdots & b & d &\cdots\cr
                &\vdots&\vdots \cr}\right]
\qquad\qquad\hbox{it holds}\qquad
b=a+c+d
$$
\end{corollary}

\begin{proof}
$(a+b)+(c+d) = A+B=2b$, hence $a+c+d=b$.
\end{proof}

The square identity is useful in producing the entries of a Krawtchouk 
matrix: fill the top row with ``1's" the right-most column with 
the binomial coefficients. Then, apply the square identity to 
reproduce the matrix. 
\\

In summary, the identities considered above can be written as follows:
\\

{\bf Cross identities:}
%
$$
\matrix{
(i)&   \FFF n{i-1\;j} + \FFF n{i  \;j}  =  \FFF {n+1}{i\;j}    &\quad&
(ii)&  \FFF n{i  \;j} + \FFF n{i  \;j+1}= 2\FFF {n-1}{i\;j}   \cr\cr
(iii)& \FFF n{i \;j}  - \FFF n{i-1\;j}  =  \FFF {n+1}{i\;j+1}  &\quad&
(iv)&  \FFF n{i  \;j} - \FFF n{i  \;j+1}= 2\FFF {n-1}{i-1\;j}\cr}
$$


{\bf Square identity:}
$$
\FFF n {ij} =  \FFF n{i-1\;j} + \FFF n{i-1\; j+1}+  \FFF n{i\;j+1}
$$

%
%{\bf Corollary:} Big vertical triangle identity (minimal case)
%$$
%-\FFF n {ij} + \FFF n{i+2\; j} = \FFF {n+2}{i+2\; j}
%$$




%===================  3   ===========================================
\section{Krawtchouk matrices from Hadamard matrices}
\label{sec:result}

The identities allow to define Krawtchouk matrices recurrently, 
with the use of tensor product.  One can ``blow up"
a (symmetric) matrix  $S^{(n)}$ by Kronecker product with the elementary
matrix
\begin{equation}
\label{eq:had}
          H=\left[\matrix{ 1& 1\cr
                           1& -1\cr}\right]
\end{equation}
%
and then contract appropriately the resulting matrix to get $S^{(n)}$.

\begin{definition}\rm
Define a square contraction $r(M)$ of an $n\times n$ matrix $M$
as an $k\times k$ matrix
\begin{equation}
\label{eq:zdzichu}
(rM)_{ij} = \sum_{a=2i-2,\; 2i-1 \atop b=2j-2,\; 2j-1}
M_{ab}
\end{equation}
where the entries of $M_{ab}$ outside of the range $(1,\ldots, n)$ 
are taken as zero.
\end{definition}

%
\begin{theorem}
Symmetric Krawtchouk matrices are:
$$
S^{(n+1)} = r ( S^{(n)}\otimes H )
$$
\end{theorem}

\begin{corollary}
Regular Krawtchouk matrices are:
$$
\Phi^{(n+1)} = r ( \Phi^{(n)} B \otimes H ) B^{-1}
$$
where $B$ is the diagonal binomial matrix defined in (\ref{eq:diagB}).
\end{corollary}


{\bf Example}.  Start with symmetric Krawtchouk matrix or order 2:
$$
S^{(3)} = 
\left[ \matrix{   1 & 2 & 1 \cr
                  2 & 0 & -2 \cr
                  1 & -2& 1  \cr} \right]
$$
Get the tensor product with $H$:
$$
S^{(3)}\otimes H = 
\left[ \matrix{   1&1 & 2&2 & 1&1 \cr
                  1&-1 & 2&-2 & 1&-1 \cr
                  2&2 & 0&0 & -2&-2 \cr
                  2&-2 & 0&0 & -2&2 \cr
                  1&1 & -2&-2& 1&1  \cr
                  1&-1 & -2&2& 1&-1  \cr} \right]
$$
surround with zeros and contract:                
$$
r (S^{(3)}\otimes H) = 
\left[ \matrix{ 0& 0& 0& 0& 0& 0&0 &0 \cr 
                0& 1&1 & 2&2 & 1&1 &0 \cr
                0& 1&-1 & 2&-2 & 1&-1&0  \cr
                0& 2&2 & 0&0 & -2&-2 &0 \cr
                0& 2&-2 & 0&0 & -2&2 &0  \cr
                0& 1&1 & -2&-2& 1&1  &0 \cr
                0& 1&-1 & -2&2& 1&-1 &0  \cr
                0& 0& 0& 0& 0& 0&0 &0 \cr } \right]
%
=
\left[ \matrix{ 1& 3&  3& 1 \cr 
                3& 3& -3& -3 \cr
                3& 3& -3& -3 \cr
                1&-3&  3& -1 \cr} \right]
$$
~
\\

These results suggest that Krawtchouk matrices can be obtained somehow from 
the pure multiple tensor powers of the elementary matrix $H$.
Such powers define a family of Hadamard matrices,
namely Sylvester-Hadamard matrices \ref{YH}.




\begin{definition}
Special Hadamard matrices are tensor products (in the matrix 
Kronecker sense) of the fundamental matrix $H$: 
%
\begin{equation}\label{eq:HH}
  H^{(n)} = 
      \underbrace{H\otimes H \ldots H}_{n\hbox{ times}} 
\end{equation}
\end{definition}



The problem is that the powers --- understood as Kronecker matrix 
products --- ``scatter" the columns and rows that would have to be 
summed up. We need to identify appropriate sets of indices to be 
summed in the contraction of the Hadamard matrices. 
\\

Define the $n$-th {\bf binary shuffling function} 
as a function
$$
         f^{(n)} : {\bf Z} \to {\bf Z}
$$
as measuring the ``binary weight" of integer numbers.
That is, if $[n_2]_i$ represents the $i-th$ digit of number $n$ 
expressed in the binary notation, then
$$
              f(x) = \sum_i [x_2]_i
$$ 
(just the number of ``ones" needed to write $x$ in the binary 
representation). Notice that, set-theoretically,  
$$       
       f(\{0,1,\ldots, 2^n-1\}) =    \{0, 1,\ldots, n\} 
$$
Here are the first 16 values of $f$ listed for the 
argument running from 0 through $2^4-1=15$:
$$
0 \quad
1 \quad
1 \quad
2 \quad
1 \quad
2 \quad
2 \quad
3 \quad
1 \quad
2 \quad
2 \quad
3 \quad
2 \quad
3 \quad
3 \quad
4 \quad
$$

\begin{remark}
\label{r:shuffle}
The shuffling function can be defined recurrently as
$$
f(0)=0 \\
f(2^n + k) = f(k)+1
$$
for $0\leq k < 2^n$.
One can thus create the sequence of the values of the shuffling function 
by starting with $0$ and then recurrently by appending a string with 
its copy with values increased by $1$:
$$
0
\ \rightarrow\
01
\ \rightarrow\
0112
\ \rightarrow\
01121223
\ \rightarrow\ 
\hbox{etc.} 
$$

\end{remark}


Now we are ready to state the result;

\begin{theorem}
Symmetric Krawtchouk matrices are reductions of Hadamard matrices:
$$
S^{(n)}_{ij} = \sum_{a\in f^{-1}(i) \atop b\in f^{-1}(j)} H^{(n)}_{ab}
             = \sum_{a\in f^{-1}(i) \atop b\in f^{-1}(j)} (\otimes^n H)_{ab}
$$
\end{theorem}




{\bf Example:}
Let us see the transformation for case $H^{(4)} \to S^{(4)}$.
First, for the typographical reasons, in the following the 
Hadamard matrices will be presented with bullets instead of numbers:
$\bullet$ for 1 and $\circ$ for -1.
Thus for instance:


\def\b{\circ}
\def\a{\bullet}
$$
H^{(1)}=\left[\matrix{
\a &\a  \cr
\a &\b  \cr 
}\right]
%
\quad
%
H^{(2)}=\left[\matrix{
\a &\a &\a &\a  \cr
\a &\b &\a &\b   \cr 
\a &\a &\b &\b   \cr
\a &\b &\b &\a   \cr
}\right]
%
\quad
%
H^{(3)}=\left[\matrix{
\a &\a &\a &\a  &\a &\a &\a &\a    \cr
\a &\b &\a &\b  &\a &\b &\a &\b    \cr
\a &\a &\b &\b  &\a &\a &\b &\b    \cr
\a &\b &\b &\a  &\a &\b &\b &\a    \cr
%
\a &\a &\a &\a  &\b &\b &\b &\b    \cr \a &\b &\a &\b  &\b &\a &\b 
&\a    \cr \a &\a &\b &\b  &\b &\b &\a &\a    \cr \a &\b &\b &\a  
&\b &\a &\a &\b    \cr }\right] 
$$ 
etc. 
Now, applying the binary 
shuffling function to $H^{(4)}$, mark the rows and columns with 
the aimed indices: 

$$
\bordermatrix{
  &0  &1  &1  &2   &1  &2  &2  &3   &1  &2  &2  &3   &2  &3  &3  &4   \cr 
0 &\a &\a &\a &\a  &\a &\a &\a &\a  &\a &\a &\a &\a  &\a &\a &\a &\a  \cr
1 &\a &\b &\a &\b  &\a &\b &\a &\b  &\a &\b &\a &\b  &\a &\b &\a &\b  \cr 
1 &\a &\a &\b &\b  &\a &\a &\b &\b  &\a &\a &\b &\b  &\a &\a &\b &\b  \cr
2 &\a &\b &\b &\a  &\a &\b &\b &\a  &\a &\b &\b &\a  &\a &\b &\b &\a  \cr
%
1 &\a &\a &\a &\a  &\b &\b &\b &\b  &\a &\a &\a &\a  &\b &\b &\b &\b  \cr
2 &\a &\b &\a &\b  &\b &\a &\b &\a  &\a &\b &\a &\b  &\b &\a &\b &\a  \cr
2 &\a &\a &\b &\b  &\b &\b &\a &\a  &\a &\a &\b &\b  &\b &\b &\a &\a  \cr
3 &\a &\b &\b &\a  &\b &\a &\a &\b  &\a &\b &\b &\a  &\b &\a &\a &\b  \cr
%
1 &\a &\a &\a &\a  &\a &\a &\a &\a  &\b &\b &\b &\b  &\b &\b &\b &\b  \cr
2 &\a &\b &\a &\b  &\a &\b &\a &\b  &\b &\a &\b &\a  &\b &\a &\b &\a  \cr
2 &\a &\a &\b &\b  &\a &\a &\b &\b  &\b &\b &\a &\a  &\b &\b &\a &\a  \cr
3 &\a &\b &\b &\a  &\a &\b &\b &\a  &\b &\a &\a &\b  &\b &\a &\a &\b  \cr
%
2 &\a &\a &\a &\a  &\b &\b &\b &\b  &\b &\b &\b &\b  &\a &\a &\a &\a  \cr
3 &\a &\b &\a &\b  &\b &\a &\b &\a  &\b &\a &\b &\a  &\a &\b &\a &\b  \cr
3 &\a &\a &\b &\b  &\b &\b &\a &\a  &\b &\b &\a &\a  &\a &\a &\b &\b  \cr
4 &\a &\b &\b &\a  &\b &\a &\a &\b  &\b &\a &\a &\b  &\a &\b &\b &\a  \cr
}
$$
which indeed, when summed up, gives the symmetric Krawtchouk
matrix $S^{(4)}$.  
Entries that produce $S_{33}$ are marked for exemplification.

%%%%%%%%%%%%%%%%%%%%% extra copy %%%%%%%%%%%%%%%%%%%%%%%%%%%%%%%%%
%$$
%H^{(4)}=\left[\matrix{
%\a &\a &\a &\a  &\a &\a &\a &\a  &\a &\a &\a &\a  &\a &\a &\a &\a  \cr
%\a &\b &\a &\b  &\a &\b &\a &\b  &\a &\b &\a &\b  &\a &\b &\a &\b  \cr 
%\a &\a &\b &\b  &\a &\a &\b &\b  &\a &\a &\b &\b  &\a &\a &\b &\b  \cr
%\a &\b &\b &\a  &\a &\b &\b &\a  &\a &\b &\b &\a  &\a &\b &\b &\a  \cr
%%
%\a &\a &\a &\a  &\b &\b &\b &\b  &\a &\a &\a &\a  &\b &\b &\b &\b  \cr
%\a &\b &\a &\b  &\b &\a &\b &\a  &\a &\b &\a &\b  &\b &\a &\b &\a  \cr
%\a &\a &\b &\b  &\b &\b &\a &\a  &\a &\a &\b &\b  &\b &\b &\a &\a  \cr
%\a &\b &\b &\a  &\b &\a &\a &\b  &\a &\b &\b &\a  &\b &\a &\a &\b  \cr
%%
%\a &\a &\a &\a  &\a &\a &\a &\a  &\b &\b &\b &\b  &\b &\b &\b &\b  \cr
%\a &\b &\a &\b  &\a &\b &\a &\b  &\b &\a &\b &\a  &\b &\a &\b &\a  \cr
%\a &\a &\b &\b  &\a &\a &\b &\b  &\b &\b &\a &\a  &\b &\b &\a &\a  \cr
%\a &\b &\b &\a  &\a &\b &\b &\a  &\b &\a &\a &\b  &\b &\a &\a &\b  \cr
%%
%\a &\a &\a &\a  &\b &\b &\b &\b  &\b &\b &\b &\b  &\a &\a &\a &\a  \cr
%\a &\b &\a &\b  &\b &\a &\b &\a  &\b &\a &\b &\a  &\a &\b &\a &\b  \cr
%\a &\a &\b &\b  &\b &\b &\a &\a  &\b &\b &\a &\a  &\a &\a &\b &\b  \cr
%\a &\b &\b &\a  &\b &\a &\a &\b  &\b &\a &\a &\b  &\a &\b &\b &\a  \cr
%}\right]
%$$
%%%%%%%%%%%%%%%%%%%%%%%%%%%%%%%%%%%%%%%%%%%%%%%%%%%%%%%%%%%%%%%%%%






%============================ 5 ==========================
\section{Krawtchouk matrices and classical random walk}
\label{sec:classical}

In this section we will give a probabilistic meaning to
the Krawtchouk matrices and some of their properties.\\

Let $\xi_i$ be independent symmetric Bernoulli random variables taking
values $\pm1$. Let $X_N=\xi_1+\cdots+\xi_N$ be the associated random
walk starting from $0$. Now observe that the generating function of
the elementary symmetric functions in the $\xi_i$ is a martingale, in
fact a discrete exponential martingale:
$$
   M_N = \prod_{i=1}^N(1+v\xi_i)=\sum_k v^k 
           \alpha_k(\xi_1,\ldots,\xi_N)
$$
where $\alpha_k$ denotes the $k^{\mathrm{th}}$ elementary symmetric
function. 
The martingale property is immediate since each $\xi_i$ has mean $0$.
Suppose that at time $N$, the number of the $\xi_i$ that are equal to $-1$
is $j_N$, with the rest equal to $+1$. Then  $j_N= (N-X_N)/2$ 
and $M_N$ can be expressed solely in terms of $N$ and $X_N$, or, 
equivalently, of $N$ and $j_N$
$$
       M_N = (1+v)^{N-j_N}(1-v)^{j_N} 
           = (1+v)^{(N+X_N)/2}(1-v)^{(N-X_N)/2}  
$$
>From the generating function for the Krawtchouk
matrices, equation (\ref{eq:genkraw}), follows 
$$
    M_N = \sum_iv^i\Phi^{(N)}_{i,j_N}
$$
so that as functions on the Bernoulli space, each sequence of random
variables $\Phi^{(N)}_{i,j_N}$ is a martingale. \\

Now we can interpret two basic recurrences of Section 2. For a fixed
column of $\Phi^{(N)}$, to get the corresponding column in 
$\Phi^{(N+1)}$ one uses the Pascal's triangle recurrence:
$$
  \FFF N{i-1\;j} + \FFF N{i  \;j}  =  \FFF {N+1}{i\;j}
$$
This follows in the probabilistic setting by writing
$M_{N+1}=(1+v\xi_N)M_N$ and remarking that for $j$ to remain constant,
$\xi_N$ must take the value $+1$, with consequences as in the purely
algebraic context. The martingale property is more interesting in the
present context. We have
$$
\FFF N{ij_N}= E(\FFF N{ij_{N+1}}|\xi_1,\ldots,\xi_N)
             =\frac12\,\left(\FFF{N+1}{i\,j_N+1}+\FFF{N+1}{ij_N}\right)
$$
since half the time $\xi_{N+1}$ is $-1$, increasing $j_N$ by 1, and
half the time $j_N$ is unchanged. Thus, writing $j$ for $j_N$,
$$
      \FFF N{ij}=\frac12\,\left(\FFF{N+1}{i\,j+1}+\FFF{N+1}{ij}\right)
$$

The interpretation of Krawtchouk polynomials as elementary symmetric
functions on the Bernoulli space allows one to derive their 
analytic properties using probabilistic methods with corresponding
interpretations. 






%========================= 6 ======================
\section{Krawtchouk matrices and quantum random walk}
\label{sec:quantum}

In quantum probability, random variables are modelled by self-adjoint
operators on Hilbert space and independence by tensor products. 
We can model symmetric Bernoulli walk as follows. 
Consider a 2-dimensional Hilbert space $V={\bf R}^2$
and two special $2\times2$ operators (matrices) 
$$ 
         F=\pmatrix{0&1\cr 1&0\cr},    \qquad
         G=\pmatrix{1&0\cr 0&-1\cr}
$$
satisfying $F^2=G^2=I$, the $2\times2$ identity.
Recall the fundamental Krawtchouk/Hadamard matrix $H$ 
(see Equation (\ref{eq:krav2}) and (\ref{eq:had})), which we shall now view as
$$
H=F+G = \left(\matrix{ 1 &  1 \cr
                       1 & -1 \cr}\right)
$$
One can readily check that
\begin{equation}
\label{eq:H}
         FH =F(F+G) =(F+G)G =HG
\end{equation}
%
(use $F^2=G^2=I$). This, of course, can be viewed as 
the spectral decomposition of $F$ and we can interpret the Hadamard
matrix as diagonalizing $F$. \\


\begin{remark}
Note that the  exponentiated  operator
$$
    \exp(zF) = \pmatrix{\cosh z&\sinh z\cr \sinh z&\cosh z\cr}
$$
has the expectation value in the state $e_0$ equal 
$$
\langle e_0,\exp(zF)e_0\rangle = \cosh z
$$
which coincides with the moment generating function for the symmetric 
Ber\-noul\-li random variable taking values $\pm 1$. 
This implies that indeed we are dealing with the (quantum)
generalization of the classical model.
\end{remark}


~

\begin{remark} {\bf (Quantum computer interpretation).}
\rm 
We can view $V=\hbox{span}\;\{e_0,\,e_1\}$ as a Hilbert space 
representing a simple two-state quantum system, e.g., 
a particle with possible spin-up ($e_0$) and spin-down ($e_1$) states, 
respectively.
In the context of ``quantum computing," these elementary states represent 
``qubits."  Operator $F$ represents a ``spin flip" and operator $G$ 
a  ``phase change,"  both are thus basic qubit operations. 
[Geometrically, they are reflections in $V$ through lines 
that form 45 degrees.  The Hadamard matrix $M$ is essentially 
a "rotation by 45 degrees" (up to composition with reflection and scaling) 
and therefore ties them via (\ref{eq:H}).] 
Now, to perform quantum computing, one needs an assembly of $N$ such
independent systems --- this constitutes a simple quantum computer.  
The Hilbert space of computer states is represented 
by the $N$-th tensor product of the original space $V$,
that is, by the $2^N$-dimensional Hilbert space $V^{\otimes N}$. 
This motivates our further considerations. 
\end{remark}



Define the following simple operators
\begin{eqnarray*}
\xi_1 &=& F\otimes I \otimes\cdots\otimes I          \\
\xi_2 &=& I\otimes F\otimes I \otimes\cdots\otimes I \\
\vdots &=& \vdots                                    \\
\xi_N &=& I\otimes I \otimes\cdots\otimes F
\end{eqnarray*}
%
Each $\xi_i$ describes a "flip" at $i$-th position.
(These are our ``quantum" equivalents of the random walk variables 
from section \ref{sec:classical}.)
Performing them independently would result in ``running" 
a classical $N$-bit computer.
We shall however consider the superposition of these independent 
individual actions,  described by $X_F=\xi_1+\cdots+\xi_N$. 
Analogously, we define:
%
\begin{eqnarray*}
\sigma_1 &=& G\otimes I \otimes\cdots\otimes I\\
\sigma_2 &=& I\otimes G\otimes I \otimes\cdots\otimes I\\
\vdots &=& \vdots \\
\sigma_N &=& I\otimes I \otimes\cdots\otimes G
\end{eqnarray*}
%
with $X_G=\sigma_1+\cdots+\sigma_N$. 
We also extend $H$ to the $N$-fold tensor product, 
setting $H_N=H^{\otimes N}$. 
\\


Now for illustration, consider a simple calculation for $N=3$:
%
\begin{eqnarray*}
\xi_1H_3 &=& (F\otimes I\otimes I)(H\otimes H\otimes H)\\
         &=& (H\otimes H\otimes H)(G\otimes I\otimes I)
              = H_3\sigma_3
\end{eqnarray*}
%
where relation $FH=HG$ was used. 
This clearly generalizes to $\xi_kH_N=H_N\sigma_k$ and 
--- by summing over $k$ --- yields an important relation:
$$
            X_F H_N = H_N X_G
$$
Since products are preserved under projection onto 
the symmetric tensor space (see Appendix 2), we get
$$
       \overline{X}_F\overline{H}_N = \overline{H}_N \overline{X}_G
$$
We know how to calculate $\overline{H}_N$ from the action of $H$ on
polynomials in degree $N$. 

\begin{proposition} For each $N>0$, symmetric projection of Hadamard 
matrices leads to 
$$ 
    \overline{H}_{ij} = \FFF{N}{ji}
$$
That is, $\overline{H}_N$ is the transpose of the Krawtchouk matrix
$\FFF{N}{}$.
\end{proposition}

\begin{proof}
Writing $(x,y)$ for $(x_0,x_1)$, we have in degree $N$ for the 
$k^{\mathrm{th}}$ component:
$$
        (x+y)^{N-k}(x-y)^k=\sum_l \overline{H}_{kl} x^{N-l}y^l
$$
Scaling out $x^N$ and replacing $v=y/x$ yields
the generating function for the Krawtchouk matrices with the
coefficient of $v^l$ equal to $\FFF{N}{lk}$. Thus the result.
\end{proof}

Now consider the generating function for the elementary symmetric
functions in the quantum variables $\xi_j$. This is the $N$-fold
tensor power
$$ 
     {\cal F}_N (t) = (I+tF)^{\otimes N}
                    = I^{\otimes N}+t\,X_N+\cdots
$$
noting that the coefficient of $t$ is $X_F$. 
Similarly, define
$$
     {\cal G}_N(t) = (I+tG)^{\otimes N}
                   = I+tX_G+\cdots
$$
  From $(I+tF)H=H(I+tG)$ we have
$$
   {\cal F}_NH_N  =  H_N{\cal G}_N        
   \qquad\hbox{and}\qquad 
   \overline{\cal F}_N \overline{H}_N  
 = \overline{H}_N \overline{\cal G}_N
$$
The difficulty is to calculate the action on the symmetric tensors for
operators, such as $X_N$, that are not pure tensor powers. However,
from ${\cal F}_N(t)$ and ${\cal G}_N(t)$ we can 
recover $X_F$ and $X_G$ via
$$
          X_F=\frac{d}{dt}\biggm|_{t=0}(I+tF)^{\otimes N},
     \qquad
          X_G=\frac{d}{dt}\biggm|_{t=0}(I+tG)^{\otimes N}
$$
with corresponding relations for the barred operators.
Calculating on polynomials yields the desired results as follows.
$$ 
          I+tF = \pmatrix{1&t\cr t&1\cr},\qquad
          I+tG = \pmatrix{1+t&0\cr 0&1-t\cr}
$$
In degree $N$, using $x$ and $y$ as variables, we get the 
$k^{\mathrm{th}}$ component for $\overline{X}_F$ and $\overline{X}_G$
via
\begin{eqnarray*}
    \frac{d}{dt}\biggm|_{t=0}(x+ty)^{N-k}(tx+y)^k
                     &=& (N-k)\dir k+1. +k\dirs k-1. \\
    \frac{d}{dt}\biggm|_{t=0}(1+t)^{N-k}(1-t)^k\dirs k.
                     &=& (N-2k)\dirs k. 
\end{eqnarray*}
respectively.  
For illustration, calculations for $N=4$ result in

\begin{eqnarray*}
\overline{X}_F = 
         \pmatrix{0&4&0&0&0\cr 
                  1&0&3&0&0\cr 
                  0&2&0&2&0\cr 
                  0&0&3&0&1\cr
                  0&0&0&4&0\cr},&\qquad&
\overline{H}_4 = \pmatrix{1&4&6&4&1\cr
                    1&2&0&-2&-1\cr 
                    1&0&-2&0&1\cr
                    1&-2&0&2&-1\cr
                    1&-4&6&-4&1\cr} \\ 
\mathstrut\\
\overline{X}_G  = \pmatrix{4&0&0&0&0\cr
                     0&2&0&0&0\cr
                     0&0&0&0&0\cr
                     0&0&0&-2&0\cr
                     0&0&0&0&-4\cr},&&
\end{eqnarray*}

We have a simple

\begin{corollary}
The spectrum of $\overline{X}_N$ is
$N,N-2,\ldots,2-N,-N$, which coincides with the support 
of the classical random walk.
\end{corollary}


We note that the top row of $(I+tF)^{\otimes N}$ is $t^{f(k)}$, where
$f(k)$ is the binary shuffling function of Section \ref{sec:result}. 
This is seen by noting that each time one tensors with $I+tF$, 
the original top row is reproduced then it is concatenated with a replica 
of itself modified in that each entry picks up a factor of $t$ 
(compare with \ref{r:shuffle}).
Now, collapsing to the symmetric tensor space,
the top row will have entries ${N\choose k}t^k$. This
follows as well by direct calculation of the $0^{\mathrm{th}}$ component
matrix elements in degree $N$: expanding $(x+ty)^N$.



%----------------------------------------------------
\subsection{Expectation values}
To find the distributions, we must calculate expectation values. In
the present context, expectation values in
two particular states are especially interesting. Namely,
in the state $e_0$ and the normalized trace --- the uniform distribution
on the spectrum. In the $N$-fold tensor product, we want to consider
expectation values in the ``ground state" $\dir 000\ldots0.$ 
and normalized traces. Then we can go to the symmetric tensors.
\\

Since everything factors, 
one easily obtains the expectation value of $\exp(zX_N)$ 
in the ground state $\dir 000\ldots0.$ to be $(\cosh z)^N$. 
For the trace, $\tr \exp(zF)=2\cosh z$ implies
$\tr\exp(zX_F)=2^N(\cosh z)^N$ and, after normalizing, this yields
$(\cosh z)^N$.
\\

For the barred operators, we consider the symmetric trace. Here we can use
the {\bf symmetric trace theorem}, detailed in Appendix 2. It tells us
that the generating function for the symmetric traces of any operator
$A$ in the various degrees is $\det(I-tA)^{-1}$. Taking
$A=\exp(zF)$, we have

\begin{eqnarray*}
  \det(I-te^{zF})^{-1} &=& [(1-te^z)(1-te^{-z})]^{-1}\\
                       &=& (1-2t\cosh z+t^2)^{-1}
\end{eqnarray*}

The latter is the generating function for Chebyshev polynomials of the
second kind, so that the normalized symmetric trace is
$$
   (N+1)^{-1}\tr_{Sym}^N \exp(zF) = U_N(\cosh z)/(N+1)
$$
This equals as well
$$
      \frac{e^{z(N+1)} - e^{-z(N+1)}}{(e^z-e^{-z})(N+1)}
$$

And that completes our study.



\vskip.3in
{\bf Acknowledgement.} We would like to thank M. Viana for inviting us
to participate in the special session and we extend our appreciation
for all the hard work involved in organizing the session as well as
related activities.

%====================== A1 ==========================
\section*{Appendix 1: Tensor products}
\label{sec:A1}


Let $V$ be a $d$-dimensional vector space over ${\bf R}$. 
We fix an orthonormal basis $\{e_0,\ldots,e_\delta\}$
with $d=1+\delta$. Denote tensor powers of $V$ by
$V^{\otimes N}$, so that $V^{\otimes 2}=V\otimes V$, etc. 
A basis for $V^{\otimes N}$ is given by all $N$-fold tensor products 
of the basis vectors $e_i$. 

$$ 
\dir n_1n_2\ldots n_N. = 
      e_{n_1}\otimes e_{n_2}\otimes \cdots\otimes e_{n_N}
$$

Note that we can label these $d^N$ basis elements by all  
numbers $0$ to $d^N-1$ and recover the tensor products 
by decomposing these numbers in numerical base $d$ 
putting leading zeros so that all extended labels are of length $N$. 
\\

Now let $\{A_i:\; 1\le i\le N\}$ be a set of $N$ linear operators on $V$. 
The linear operator on $V^{\otimes N}$, 
$A=A_1\otimes A_2\otimes\cdots\otimes A_N$ acts on a basis vector
$\dir n_1n_2\ldots n_N.$ by
$$ 
 A \dir n_1n_2\ldots n_N. =
       A_1e_{n_1}\otimes\cdots\otimes A_Ne_{n_N}
$$
This needs to be expanded and terms regrouped using bilinearity.\\

If $A$ and $B$ are two $d\times d$ matrices, the matrix corresponding
to the operator $A\otimes B$ is the Kronecker product, a
$d^2\times d^2$ matrix having the block form:
$$\pmatrix{
       a_{00}B&\ldots&a_{0\delta}B\cr
       a_{10}B&\ldots&a_{1\delta}B\cr
              \vdots&\vdots&\vdots\cr
       a_{\delta0}B&\ldots&a_{\delta\delta}B\cr}
$$
and iteratively, associating from the left, for higher-order tensor
products.
The rows and columns of the matrix of a linear operator acting on
$V^{\otimes N}$ are conveniently labelled by associating to each basic tensor
$\dir d_1d_2\ldots d_N.$ the corresponding integer label
$\sum\limits_{k=1}^N n_kd^{N-k}$, which thus provides a canonical
ordering.


%======================= A2 ============================
\section*{Appendix 2: Symmetric tensor spaces}
\label{sec:A2}

Here we review relation between tensor spaces and symmetric tensor spaces
and the role of symmetric polynomials.
\\

 
The space $V^{\otimes N}$ can be mapped onto the space of symmetric
tensors, $V^{\otimes_S N}$ by identifying basis vectors 
(in $V^{\otimes N}$) that are equivalent under all permutations. 
Alternatively, one can identify the basic tensor
$\dir d_1d_2\ldots d_N.$ with the monomial $x_{d_1}x_{d_2}\cdots x_{d_N}$
in the commuting variables $x_0,\ldots,x_\delta$. 
This equivalence class will be denoted $\dirs d_1d_2\ldots d_N.$ 
Hence we have a linear map
from tensor space into the space of symmetric polynomials, itself 
isomorphic to the space of symmetric tensors:

\begin{equation}
\label{eq:proj}
  \overline{\phantom{W}}:\quad V^{\otimes N} 
                \ \longrightarrow \ 
                \hbox{Sym}(x)  \cong V^{\otimes_S N}
\end{equation}

In the symmetric tensor space, tensor labels
need to count only `occupancy,' that is, the number of times a
basis vector of $V$ occurs in a given basic tensor of $V^{\otimes N}$. 
We indicate occupancy by a multi-index which is the exponent of the
corresponding monomial. 
The dimension of $V^{\otimes_S N}$ is thus
$$
     \hbox{dim}\; V^{\otimes_S N} = {N+d-1\choose d-1}
$$
that is, the number of monomials homogeneous of degree $N$. 
A canonical labelling is $\dirs k_1\ldots k_\delta.$, 
with $k_0$ understood to equal $N-\sum\limits_{j=1}^\delta k_j$. 
\\

Given an operator $A$ on $V$, let $A_N=A^{\otimes N}$. Then $A_N$  
induces an operator $\overline{A}_N$ on $V^{\otimes_S N}$ from the action of
$A$ on polynomials,  which we call the {\bf symmetric representation 
of $A$ in degree $N$}.  Denote the matrix elements of the action of
$\overline{A}_N$ by $\overline{A}_{mn}$. If $A$ has matrix entries
$A_{ij}$ let
$$
       y_i=\sum_i A_{ij}x^j
$$
Then the matrix elements are defined by the relation (expansion):
$$
   y_0^{m_0}\cdots y_\delta^{m_\delta}=
               \sum_n \overline{A}_{mn}
                     x_0^{n_0}\cdots x_\delta^{n_\delta}
$$
with multi-indices $m$ and $n$.
\\

Successive application of $A_1$ then $A_2$ shows that mapping 
(\ref{eq:proj}) is an algebra homomorphism, i.e., 
$$
\overline{A_1A_2}=\overline{A}_1\overline{A}_2
$$
Explicitly, in basis notation
$$ 
\overline{(A_1A_2)}_{mn}=
          \sum_r (\overline{A_1})_{mr}
                 (\overline{A_2})_{rn}
$$ 

Define the {\bf symmetric trace} in degree $N$ of $A$ as the trace of
the matrix elements of $\overline{A}_N$, i.e., the sum of the diagonal
matrix elements:
$$
     \tr_{\rm Sym}^N A=\sum_{|m|=N}\overline{A}_{mn}
$$
with $|m|$ denoting ---as usual--- the sum of the components of $m$.
Observe that if $A$ is upper-triangular, 
with eigenvalues $\lambda_1,\ldots,\lambda_d$, then the trace of
this action on the space of polynomials homogeneous of degree $N$ is 
exactly  $h_N(\lambda_1,\ldots,\lambda_d)$, the $N^{\mathrm{th}}$ 
homogeneous symmetric function in the $\lambda$'s. 
We recall a useful theorem on calculating the symmetric trace.
Since the mapping from $A$ to $\overline{A}_N$ is a homomorphism, a
similarity transformation on $A$ extends to one on $\overline{A}_N$ thus
preserving traces. Now, any matrix is similar to an upper-triangular 
one with the same eigenvalues, thus follows (\cite{Spr}):

\begin{thm}{\bf Symmetric trace theorem}
\label{thm:symmtr}
Denoting by $\mathrm{tr}_{\mathrm{Sym}}^N$ the trace of the symmetric
representation on polynomials homogeneous of degree $N$, 
$$
     \frac{1}{\det(I-tA)}=\sum_{N=0}^\infty t^N\tr_{\mathrm{Sym}}^N A
$$
\end{thm}

\begin{proof} 
With $\{\lambda_i\}$ denoting the eigenvalues of $A$,
\begin{eqnarray*}
\frac{1}{\det(I-tA)} 
           &=&\prod_i\frac{1}{1-t\lambda_i}\cr
           &=&\sum_{N=0}^\infty t^Nh_N(\lambda_1,\ldots,\lambda_d)\cr
           &=&\sum_{N=0}^\infty t^N\tr_{\mathrm{Sym}}^N A\cr
\end{eqnarray*}
as stated above.
\end{proof}

%%%%%%%%%
%Observe that with 
%$\displaystyle \det(I-tA)=1-\sum_{k\ge1} a_kt^k$, $\det(I-tA)^{-1}$ is
%the generating function for the fundamental solution to the recurrence
%$$ 
%   y_n=\sum_{k=1}^d a_k y_{n-k}
%$$
%In particular, if $A$ is a $2\times2$ matrix, with $\tr A=\tau$, 
%$\det A=\Delta$, then $\det(I-tA)^{-1}$ is the generating function for 
%the fundamental solution to the recurrence
%\begin{equation}
%  \label{eq:tau-delta}
% y_n=\tau y_{n-1}-\Delta y_{n-2} 
%\end{equation}
%%%%%%%%%%%%%%%%%%%%



If $A$ is represented by a matrix, given the matrix form 
$A_N=A^{\otimes N}$ computed as an $N$-fold Kronecker product, to reduce
to $\overline{A}_N$, we see that acting on polynomials, for a fixed row
label in the full tensor space, the column entries corresponding to
basic tensors equivalent under permutation are summed to a single
column. Then the matrix elements are chosen, one row from each
equivalence class of basic tensors. What if $d=2$? Then row and column
labels are single indices so that $\overline{A}_N$ is an 
$(N+1)\times(N+1)$ matrix with labels according to the exponent of $x_1$.
That is, the basis for polynomials in degree $N$ is given by the
monomials $x_0^{N-k}x_1^k$, $0\le k\le N$. The binary shuffling
function and contraction operations of Section \ref{s:result} 
are exactly the reduction to symmetric tensors and then computing 
the induced matrix, respectively. 





%==================================================================
\begin{thebibliography}{10}

\bibitem{B}
Bose N:
\emph{Digital filters: theory and applications},
North-Holland, 1985.


\bibitem{Del}
Delsarte P:
Bounds for restricted codes, by linear programming,
\emph{Philips Res. Reports}, {\bf 27} (1972) 272--289.

\bibitem{Del2}
Delsarte P:
Four fundamental parameters of a code and their combinatorial significance,  
\emph{Info. \&\ Control}, {\bf 23} (1973) 407--438

\bibitem{Del3}
Delsarte P:  
An algebraic approach to the associacion schemes of coding theory, 
\emph{Philips Research Reports Supplements}, No. 10 (1973)

\bibitem{Dun}
Dunkl C F:
A Krawtchouk polynomial addition theorem and  wreath products of symmetric
groups, 
\emph{Indiana Univ. Math. J.},  {\bf 25} (1976) 335--358

\bibitem{DR}
Dunkl C F and Ramirez D F:
Krawtchouk polynomials and the symmetrization of hypergraphs,
\emph{SIAM J. Math. Anal.}, {\bf 5} (1974) 351--366

\bibitem{Eag}
Eagelson G K:
A characterization theorem for positive definite sequences of the
Krawtchouk polynomials,  
\emph{Australian J. Stat}, {\bf 11} (1069) 29--38

\bibitem{FF}
Feinsilver P and Fitzgerald R: 
The spectrum of symmetric Krawtchouk matrices, 
\emph{Lin. Alg. \& Appl.}, 235, 1996, 121-139. 

\bibitem{FS}
Feinsilver P and Schott R:
Krawtchouk polynomials and finite probability theory,
\emph{Probability Measures on Groups X}, Plenum, 1991, pp. 129--135

\bibitem{FH}
Fulton W and Harris J:
\emph{Representation theory, a first course,}
Graduate texts in mathematics, {\bf 129},
Springer-Verlag, 1991.

\bibitem{Kra}
Krawtchouk M:
Sur une generalisation des polynomes d'Hermite,
\emph{Comptes Rendus}, {\bf 189} (1929) 620--622

\bibitem{Kra2}
Krawtchouk M: 
Sur la distribution des racinesdel polynomes orthogonaux,
\emph{Comptes Rendus}, {\bf 196} (1933) 739--741

\bibitem{Lev}
Levenstein V I:
Krawtchouk polynomials and universal bounds for codes and design in Hamming
spaces,
\emph{ISEE Transactions on Information Theory}, {\bf 41} 5 1995 (sep)  pp
1303--1321

\bibitem{MS}
MacWilliams F J and Sloane N J A: 
\emph{The theory of Error-Correcting Codes.}
The Netherlands, North Holland, 1977.

\bibitem{Spr}
Springer T A:
\emph{Invariant theory,}
Lecture notes in mathematics, {\bf 585},
Springer-Verlag, 1977.


\bibitem{Sze}
Szeg\"o G:
{\sl Orthogonal Polynomials}, 
Colloquium Publications, Vol. 23, New York, AMS, revised eddition 1959, p.
35--37


\bibitem{V-J}
Vera-Jones D:
Finite bivariate distributions and semi-groups of nonnegative matrices,
\emph{Q. J. Math. Oxford}, {\bf 22} 2 (1971)
247--270.


\bibitem{YH} Yarlagadda R K Rao and Hershey J E:
\emph{Hadamard matrix analysis and synthesis: with applications
to communications and signal/image processing,}
Kluwer Academic Publishers, 1997.

\end{thebibliography}



\end{document}

